\section{Discussion}
\label{sec:discussion}

\subsection{What a physics-first inverse model buys}

The central argument of this paper is not that \alibz{} is more accurate than
calibrated or learned methods; Section~\ref{sec:results} shows that the absolute,
closed-sum quantification of heterogeneous single shots is not yet reached. The
argument is that an inverse model built from Equations~\eqref{eq:emissivity},
\eqref{eq:saha} and \eqref{eq:wnnls} fails in ways that can be seen. Every failure
mode documented in Section~\ref{sec:results} was caught by a diagnostic that the
model itself produces: a basin collapse appears as a top fraction that jumps to one
while $r^2$ \emph{improves}; a confounded attribution appears as a supporting peak set
that a rival could cover at its own true-negative cap (Equation~\eqref{eq:contest});
a phantom trace element appears as a detection score $z<3$ on its strongest predicted
line (Equation~\eqref{eq:detect}); a wrong plasma state appears as a stage
disagreement $D_\mathrm{el}$ (Equation~\eqref{eq:stage-disagreement}) or an electron
density pinned to the edge of its prior box. A purely empirical regressor
\citep{clegg2009pls,li2020cnn} returns a composition in each of these cases with no
residual structure to flag it, because nothing in the model is obliged to be
consistent with Saha--Boltzmann physics. That is the sense in which the assumptions of
Section~\ref{sec:assumptions} are an asset rather than a liability: the ones that
matter are tested at run time, and the ones that are not yet tested are named.

\subsection{Where the model is limited by physics, and where by the instrument}

Two limitations are intrinsic to the measurement and no amount of software removes
them. The first is the weak information that line amplitudes carry about the plasma
state. The amplitude objective is flat in $\nel$ at fixed $T$ and nearly flat in $T$
over several thousand kelvin (Section~\ref{sec:theory-thermometer}), so the
composition of a feldspar-like spectrum, whose Si/K emissivity ratio changes by a
factor of $\sim$2.5 per \SI{1000}{\kelvin}, is only as good as the thermometer and
the density anchor. The stage-consistency cost fixes $T$ where two stages of one
element are observed, and H$\alpha$ fixes $\nel$ where hydrogen is present
\citep{gigosos2003halpha}; otherwise \alibz{} falls back to a prior and says so. This
is the same identifiability problem that limits classical CF-LIBS Boltzmann and
Saha--Boltzmann plots \citep{tognoni2010cflibs,aguilera2007multielement}, made
explicit by treating the stages as independent unknowns rather than hidden by tying
them. The second intrinsic limitation is spatial and temporal inhomogeneity
\citep{aragon2008review,cristoforetti2013columnar}: a single $(T,\nel)$ is an
effective average over a gated, expanding plasma, and on a multi-phase mineral the
same element is excited in different host phases. The stage disagreement is the
diagnostic for both, and it cannot separate them; correlative characterization of the
ablated phases is the only route to doing so.

Three further limitations belong to the instrument and are therefore removable. The
relative spectral response within and between the ultraviolet and visible segments is
uncalibrated (Assumption~A6), and the \SI{620}{\nano\meter} gain step is measured
from the continuum with a corpus fallback; on feldspar-like spectra the K/Si split is
limited by that fallback, not by the thermometer. A lamp calibration and a
branching-ratio method \citep{li2014spectralresponse,urbinamedina2021branching} are
the planned fixes. The line-spread function is not a Voigt, the vendor export grid is
a spline resample of coarser native pixels, and the pipeline now fits on the native
grid with per-segment location estimators (Section~\ref{sec:wf-fit}); a
forward-convolved, pixel-integrated line-spread function remains to be built, and
until it is, the self-absorption parameters of Equation~\eqref{eq:sa-voigt} partly
absorb instrument shape (Assumption~A9). Finally, the wavelength axis carried an
18-pixel offset that invalidated every line identification made before it was found
(Section~\ref{sec:results}); the per-spectrum registration against ambient argon
lines now bounds the residual to a fraction of a pixel, but it is a reminder that an
instrument-level error can masquerade as physics.

Atomic data are a limitation of a third kind. Transition probabilities enter
Equation~\eqref{eq:emissivity} linearly and carry accuracy classes from a few percent
to \SI{50}{\percent} \citep{kramida2023asd,fuhr2006fe}; partition functions are
summed from the line-derived level set rather than taken from critically evaluated
tables \citep{barklem2016partition}; ionization-potential depression and the
McWhirter criterion are not evaluated (Section~\ref{sec:theory-saha}). None of these
is tested at run time. They set a floor on the accuracy that any calibration-free
method can reach, and the synthetic generator samples them as source-aware
uncertainty precisely so that a trained model does not learn to trust them.

\subsection{The MW2-112 production gate}

The clearest statement of where the method stands is the pre-registered production
gate of the MW2-112 stratigraphic profile (Section~\ref{sec:datasets}). The
closed-sum CF-LIBS branch was rejected: in the 20-shot quantitative pilot only 3
positions passed the physical quality control, 4 warned and 12 failed on low or
negative pattern $r^2$, plasma parameters at bounds, or weak-shape and composition
collapse, and doubling the optimizer effort or varying the \SI{620}{\nano\meter}
fallback response changed the compositions without stable improvement.
% src: docs/mw2_112/analysis_plan.md §9 (execution record 2026-07-15)
A relative-profile branch, which asks only for within-element spatial ranks from
same-stage multiplets, was completed for all 929 spectra and validates against the
vendor's independent summaries (Spearman $\rho$ of 0.93 for Si, 0.88 for Ti, 0.71 for
Fe, 0.66 for Li and 0.60 for Mg).
% src: docs/mw2_112/analysis_plan.md §9
The contrast is instructive: the forward model and the peak-extraction chain are
good enough to rank, and the identification and plasma-state inference are the
weak links for absolute quantification on single shots of a heterogeneous,
finely stratified mudrock. The calibration campaign on pure-element standards
(Section~\ref{sec:datasets}) was begun for exactly this reason: on a target whose
composition is known, every residual is attributable to the instrument or the
physics rather than to the sample.

\subsection{What the panel taught: basins compensate, matrices must conserve}

Two lessons from Section~\ref{sec:results-ab} generalise beyond the specific changes. First, a
physically more correct ingredient can score worse on a round trip when the baseline's error is
being cancelled by a second error: the line-sum partition function landed the search on a cold
basin that compensated a composition bias, and the tabulated functions were penalised for
removing the compensation. A fixed panel with a solve at the known true state is what exposed
this; a benchmark that only reports end-to-end accuracy would have rejected the correct physics.
Second, the whole-pattern design matrix must conserve emission. Spreading one database line over
several nearby peaks is a natural way to tolerate wavelength error, but it lets one line be
counted against two fitted components, which over-predicts the elements whose lines most often
have a second component within the tolerance, and the non-negative solve then under-recovers
them; on the silicate scene that is silicon, and on real spectra it is likely to be iron, whose
forest is the densest. The principled alternative is to let each fitted component be the sum of
the database lines under it. This is a property of the inverse model, not of the plasma, and
it is invisible to goodness of fit. We expect it to account for part of the single-element
dominance failures of Section~\ref{sec:results} and for the pure-iron standard's low iron
fraction; the conserving matrix is therefore the first change of the next round, and the
candidates already scored are re-evaluated on top of it rather than on the biased solver.

\subsection{Relation to learned models}

We regard the physics-first inverse model and the learned models of
Section~\ref{sec:intro} as complementary rather than competing, in two concrete
ways. First, the forward model is the source of labelled training data. The
explicit-stage synthetic generator (Section~\ref{sec:datasets}) renders spectra with
independent stage abundances, per-stage temperatures and source-aware atomic
uncertainty onto arbitrary pixel grids, and the planned coordinate-aware transformer
is specified to be trained on it, with acceptance gates on conditional peak-profile
and export-kernel realism before training begins. Randomizing the simulator in this
way is the established route to transfer from synthetic to real data
\citep{tobin2017domainrand}, and set- and point-based attention
\citep{lee2019settransformer,zaheer2017deepsets} is the natural architecture for
irregularly sampled, multi-segment spectra. Second, the inverse model supplies the
consistency checks that a learned model lacks: a network's composition can be
forward-modelled at its implied plasma state and scored by the same residual, stage
disagreement and detectability diagnostics as a solver's. Physics-informed training
\citep{raissi2019pinn} embeds the same equations as penalties; the difference is
whether the physics is a constraint on the answer or a diagnostic of it, and we
argue for keeping both.

A more immediate statistical improvement is already identified. Comparing basins by
$r^2$ fails because a degenerate single-element basin can fit a noisy spectrum
better than the truth (Section~\ref{sec:results}); whitening the objective by a
per-peak noise model (Equation~\eqref{eq:sigma-eff}) repairs the comparison on
synthetic scenes but remains opt-in because its $\chi^2$ still misranks
temperature. A pixel-level likelihood with a parametric, positive-definite noise
covariance, replacing the peak table as the solver's input, is the fundamental fix;
its prototype showed that correlated-noise whitening on the resampled export grid
rewards sharp features and must be paired with the native grid and a noise nugget.

\subsection{Open validation gates}

Three gates remain open, in order of what they would establish. (i)
\emph{Composition invariance}: the same rock measured under argon and in air must
return the same composition while $T$ and $\nel$ differ \citep{palasti2022argon};
the metric is defined (Section~\ref{sec:metrics}) but no measured pair has yet been
staged into the repository, so it has been exercised only on a synthetic argon-ambient
scene. (ii) \emph{Pure-element standards}: the Z300 iron and vanadium studies
constrain the plasma temperature, the doublet optical depths and the wavelength
registration, but a pure metal cannot test the composition solve. (iii)
\emph{Certified reference materials}: no certified geological or glass standard
\citep{jochum2005georem,jochum2011nist} has yet been measured; the benchmark
scaffolding carries placeholders for them, and the working accuracy target is a
relative error on major elements of about \SI{20}{\percent}.
% src: DECISIONS.md 2026-11-03 entry recording the 2026-07-01 target
Until (iii) is passed, every accuracy statement about \alibz{} rests on synthetic
round trips, whose noise and atomic data are, by construction, those of the model.
